令 $U(t)=e^{-it|D|}$。取标量 Fourier 变换 $\widehat\psi(t)=\int e^{-itu}\psi(u)\,du$，由反演得
\[
 A_{\delta,r}g(x)=\frac{\delta r}{2\pi}
 \int_{\mathbb R}e^{itr}\widehat\psi(\delta rt)U(t)g(x)\,dt.
\]
关于 $1\le r\le2$ 的上确界由 $C_M\delta\int(1+\delta|t|)^{-M}|U(t)g(x)|\,dt$ 控制。在 $R\le|t|\le2R$ 上，其中 $R\ge1$ 为二进尺度，对时间变量用 Hölder 不等式并结合 \ref{ms:lemma-10-3}，得到 $L_x^p$ 界，所估计的贡献来自
\[
 C_M\delta(1+\delta R)^{-M}R^{1-1/p}R^{1/p+\sigma}\|g\|_p.
\]
对 $M>1+\sigma$，和为 $O(\delta^{-\sigma}\|g\|_p)$，可在 $R=\delta^{-1}$ 处分割求和得到。$|t|\le1$ 的贡献为 $O(\delta\|g\|_p)$，这是由 \ref{ms:lemma-10-3} 中的短时估计得到的。
